\begin{exercisechunk}
\item \textbf{Maximum likelihood need not control forward KL.}
\label[exercise]{ex:l18-likelihood-support}
\exerciseprofile{2}{1}{1}{\excalculation\quad\exinterpretation}{%
Explain why a finite-sample likelihood theorem can control Hellinger
distance without controlling forward KL.}
Let $P=\operatorname{Bernoulli}(1/2)$ and let $Q$ put probability one
on $0$. Observe $n\ge1$ independent samples $X_1,\ldots,X_n$ from $P$,
and maximize their likelihood over the realizable class $\{P,Q\}$.
\begin{enumerate}[(a)]
\item On $E=\{X_1=\cdots=X_n=0\}$, identify the unique maximum-likelihood
model and show that $\P(E)=2^{-n}$.
\item Show that $\KL(P\Vert Q)=\infty$, $h^2(P,Q)=2-\sqrt2$ using the
lecture's normalization, and $\TV(P,Q)=1/2$. Explain why part (a) prevents a general
claim that empirical maximum likelihood always has finite forward KL.
\end{enumerate}
\end{exercisechunk}
